% META: {"id":"1SPE-DERIVATION-GLOBAL-RE-C3","chapitre":"1SPE-DERIVATION-GLOBAL","type_objet":"remediation","capacites_codes":["C3"],"status":"generated"}

%---------------------------------------------------------------
% FICHE DE REMÉDIATION — C3 : Signe de la dérivée et tableau de variations
%---------------------------------------------------------------

\begin{exercice}{1SPE-DERGLOBAL-RE-C3-EX1}{1}{9}
Soit $f(x) = x^2 - 6x + 5$ définie sur $\mathbb{R}$.
\begin{enumerate}
  \item Calculer $f'(x)$.
  \item Résoudre $f'(x) = 0$ et dresser le tableau de signes de $f'$.
  \item En déduire le tableau de variations de $f$. Préciser les valeurs de $f$ aux points remarquables.
\end{enumerate}
\end{exercice}

\coupDePouce{1}{$f'(x) = 2x - 6$. Résoudre $2x - 6 = 0$, puis étudier le signe de $2x - 6$ en fonction de la position de $x$ par rapport à la racine.}

% BEGIN-VERIFY
% from sympy import symbols, diff, solve, simplify, Rational
% x = symbols('x')
% f = x**2 - 6*x + 5
% fp = diff(f, x)
% assert fp == 2*x - 6
% roots = solve(fp, x)
% assert roots == [3]
% assert fp.subs(x, 0) == -6   # negative before root
% assert fp.subs(x, 4) == 2    # positive after root
% assert f.subs(x, 3) == -4    # minimum value
% END-VERIFY

\begin{corrige}{1SPE-DERGLOBAL-RE-C3-EX1}

\textbf{Question 1.}
$f'(x) = 2x - 6.$

\textbf{Question 2.}
$f'(x) = 0 \Leftrightarrow 2x - 6 = 0 \Leftrightarrow x = 3.$

\begin{center}
\begin{tabular}{|c|ccccc|}
\hline
$x$ & $-\infty$ & & $3$ & & $+\infty$ \\
\hline
$f'(x)$ & $-$ & $-$ & $0$ & $+$ & $+$ \\
\hline
\end{tabular}
\end{center}

\commentaireMarge{$f'(x) < 0$ pour $x < 3$ (on prend $x=0$ : $f'(0) = -6 < 0$) et $f'(x) > 0$ pour $x > 3$.}

\textbf{Question 3.}
$f(3) = 9 - 18 + 5 = -4.$

\begin{center}
\begin{tabular}{|c|ccccc|}
\hline
$x$ & $-\infty$ & & $3$ & & $+\infty$ \\
\hline
$f(x)$ & $+\infty$ & $\searrow$ & $-4$ & $\nearrow$ & $+\infty$ \\
\hline
\end{tabular}
\end{center}

$f$ admet un minimum global de $-4$ en $x = 3$.

\end{corrige}

%---------------------------------------------------------------

\begin{exercice}{1SPE-DERGLOBAL-RE-C3-EX2}{1}{9}
Soit $g(x) = -2x^3 + 6x$ définie sur $\mathbb{R}$.
\begin{enumerate}
  \item Calculer $g'(x)$ et le factoriser.
  \item Dresser le tableau de signes de $g'(x)$.
  \item En déduire le tableau de variations de $g$.
\end{enumerate}
\end{exercice}

% BEGIN-VERIFY
% from sympy import symbols, diff, factor, solve, simplify, simplify
% x = symbols('x')
% g = -2*x**3 + 6*x
% gp = diff(g, x)
% assert simplify(gp - (-6*x**2)) == 0 + 6
% assert simplify(factor(gp) - (-6*(x-1)*(x+1))) == 0
% roots = solve(gp, x)
% assert set(roots) == {-1, 1}
% assert gp.subs(x, 0) == 6    # positive between roots
% assert gp.subs(x, 2) == -18  # negative after x=1
% assert g.subs(x, 1) == 4     # local max
% assert g.subs(x, -1) == -4   # local min
% END-VERIFY

\begin{corrige}{1SPE-DERGLOBAL-RE-C3-EX2}

\textbf{Question 1.}
$g'(x) = -6x^2 + 6 = -6(x^2 - 1) = -6(x-1)(x+1).$

\textbf{Question 2.}

\begin{center}
\begin{tabular}{|c|ccccccc|}
\hline
$x$ & $-\infty$ & $-1$ & & $1$ & & & $+\infty$ \\
\hline
$-6$ & $-$ & $-$ & $-$ & $-$ & $-$ & $-$ & $-$ \\
$x-1$ & $-$ & $-$ & $-$ & $0$ & $+$ & $+$ & $+$ \\
$x+1$ & $-$ & $0$ & $+$ & $+$ & $+$ & $+$ & $+$ \\
$g'(x)$ & $-$ & $0$ & $+$ & $0$ & $-$ & $-$ & $-$ \\
\hline
\end{tabular}
\end{center}

\textbf{Question 3.}
$g(-1) = 2 - 6 = -4$ et $g(1) = -2 + 6 = 4.$

\begin{center}
\begin{tabular}{|c|ccccccc|}
\hline
$x$ & $-\infty$ & & $-1$ & & $1$ & & $+\infty$ \\
\hline
$g(x)$ & $+\infty$ & $\searrow$ & $-4$ & $\nearrow$ & $4$ & $\searrow$ & $-\infty$ \\
\hline
\end{tabular}
\end{center}

\end{corrige}

%---------------------------------------------------------------

\begin{exercice}{1SPE-DERGLOBAL-RE-C3-EX3}{1}{9}
Soit $h(x) = \dfrac{x^2 + 1}{x}$ pour $x > 0$.
\begin{enumerate}
  \item Simplifier l'expression de $h(x)$ en séparant le quotient.
  \item Calculer $h'(x)$.
  \item Étudier le signe de $h'(x)$ sur $]0\,;\,+\infty[$ et dresser le tableau de variations de $h$.
\end{enumerate}
\end{exercice}

% BEGIN-VERIFY
% from sympy import symbols, diff, simplify, solve, Rational
% x = symbols('x', positive=True)
% h = (x**2 + 1)/x
% h2 = x + 1/x
% assert simplify(h - h2) == 0
% hp = diff(h2, x)
% assert hp == 1 - 1/x**2
% # h'(x) = 0 when x=1 (x>0)
% assert simplify(hp.subs(x, 1)) == 0
% assert hp.subs(x, 2) > 0
% assert hp.subs(x, x/2).subs(x, 1) < 0  # at x=0.5
% END-VERIFY

\begin{corrige}{1SPE-DERGLOBAL-RE-C3-EX3}

\textbf{Question 1.}
$h(x) = \dfrac{x^2 + 1}{x} = \dfrac{x^2}{x} + \dfrac{1}{x} = x + \dfrac{1}{x}.$

\textbf{Question 2.}
$h'(x) = 1 - \dfrac{1}{x^2} = \dfrac{x^2 - 1}{x^2}.$

\commentaireMarge{$x^2 > 0$ pour $x > 0$, donc le signe de $h'$ est celui du numérateur $x^2-1$.}

\textbf{Question 3.}
Pour $x > 0$ : $h'(x) = 0 \Leftrightarrow x^2 - 1 = 0 \Leftrightarrow x = 1$ (racine positive).

\begin{center}
\begin{tabular}{|c|ccccc|}
\hline
$x$ & $0$ & & $1$ & & $+\infty$ \\
\hline
$h'(x)$ & $-$ & $-$ & $0$ & $+$ & $+$ \\
\hline
$h(x)$ & $+\infty$ & $\searrow$ & $2$ & $\nearrow$ & $+\infty$ \\
\hline
\end{tabular}
\end{center}

$h(1) = 1 + 1 = 2$ est le minimum de $h$ sur $]0\,;\,+\infty[$.

\end{corrige}
